\chapter{The project as files}\label{ch:project-fs}
\begin{goals}
\textbf{In this chapter} you read and change a project through \file{/project}, where every dataset,
diagram, layer and figure is a file.
\end{goals}

The console starts in \file{/project}. It is not a copy of the project: it \emph{is} the project, shown
as files. Read a file and you read the project as it is now; write one and you change the project,
through the same rules as the menus, and the change can be undone like any other.

\section{The tree}\label{sec:project-tree}

\begin{center}\small
\begin{tabularx}{\linewidth}{@{}>{\raggedright}p{0.5\linewidth}L@{}}\toprule
\sffamily\bfseries path & \sffamily\bfseries what it is \\\midrule
\file{project.json} & the project in brief \\
\file{data/tocher.csv} & the rows \\
\file{data/tocher.json} & kind, columns, roles, lineage, where it is drawn \\
\file{diagrams/Net 1/diagram.json} & the diagram's settings \\
\file{diagrams/Net 1/net.svg} & the diagram as drawn \\
\file{diagrams/Net 1/layers/tocher.csv} & a link: \file{-> ../../../data/tocher.csv} \\
\file{diagrams/Net 1/layers/tocher.json} & visible, style, what it draws \\
\file{figures/<name>.compo.json} & a figure page; \file{<name>.pdf} is the page as a PDF \\
\file{selection/tocher.rows} & the selected rows, numbered from 1 \\
\file{selection/tocher.csv} & the selected rows themselves \\
\file{field/<name>/} & an original field log, as it came \\
\bottomrule\end{tabularx}
\end{center}

A layer is a \term[symbolic link]{symbolic link} to its dataset. That is the \enquote{data once, layers
per diagram} rule of \cref{sec:data-once} made visible: \code{ls -l} shows every layer pointing at the
one file that holds the rows.

\section{Reading}

Everything reads with the usual tools. \file{net.svg} is the diagram as vector graphics, ready to copy
out; \file{.pdf} files render the figure when you read them.

\begin{lstlisting}
head -5 data/tocher.csv
cat diagrams/Net\ 1/layers/tocher.json
grep -c . selection/tocher.rows
\end{lstlisting}

\section{Writing}\label{sec:project-write}

A write is an action on the project:

\begin{description}
\item[A dataset's CSV.] Writing \file{data/<name>.csv} replaces its rows; every layer redraws. A new
name in \file{data/} is a new dataset, read the way \menu{File>Open} reads a file.
\item[A settings file.] Writing \file{diagram.json} changes the diagram (its name, projection,
rotation); writing a layer's \file{.json} shows, hides or restyles what it draws.
\item[A selection.] Writing row numbers to \file{selection/<name>.rows} selects those rows; \code{rm}
clears the selection.
\item[A figure.] Writing \file{figures/<name>.compo.json} replaces the page; a new name is a new figure.
\end{description}

And the directory operations mean what they look like:

\begin{description}
\item[\code{mkdir diagrams/Joints}] makes a new net called \ui{Joints}.
\item[\code{ln -s ../../../data/tocher.csv diagrams/Joints/layers/}] draws \code{tocher} on it. The
link's name is the layer's name.
\item[\code{mv}] renames a dataset, a diagram or a figure; moving a layer's file into another diagram's
\file{layers/} moves the layer there.
\item[\code{rm -r diagrams/Joints}] removes the net.
\end{description}

\begin{why}
The rules don't loosen because you are typing. \code{rm data/tocher.csv} refuses while \code{tocher}
is drawn somewhere, as deleting it from the tree asks first; remove its layers, then the data. Files
that are computed (\file{net.svg}, \file{.pdf}, \file{selection/<name>.csv}) are read-only and say so.
\end{why}

\begin{tryit}
\begin{enumerate}
\item With the Tocher sample open, type \code{mkdir diagrams/Second}. A new tab, \ui{Second}, appears.
\item Type \code{ln -s ../../../data/tocher.csv diagrams/Second/layers/}. The Tocher poles are drawn on
it, and the tree marks the dataset as drawn twice.
\item Type \code{echo 1 2 3 > selection/tocher.rows}. Three poles turn violet on both nets.
\item Type \code{undo} twice. The layer goes, then the net. The selection goes back to what it was
before them: choosing rows is never an undo step of its own, but each step puts back the selection it
found (\cref{sec:undo}).
\end{enumerate}
\end{tryit}

\section{Field originals}

\file{field/} holds the field logs brought into the project, each in a directory of its own: the CSV
as the phone wrote it, its \file{.sig}, and a \file{photos/} directory. They are read-only. An original
goes whole (\code{rm -r field/<name>}) or not at all, and nothing can be written into one, because its
signature is over exactly those bytes. \cref{sec:field-logs} explains why.
